% GENERATED FILE. Edit canonical lessons, then run npm run generate:books.
% canonical-source-hash: fnv1a64:0643114246dbca23
% canonical-lessons: ES-C26-agua-vino, ES-C26-pan

\chapter{Water, Wine, and Bread}
\label{ch:spanish-water-wine-bread}

This chapter is generated from the canonical micro-lessons used by Language Ladder.
Each section stays independently resumable and preserves the authored prerequisite order.

\section[el agua, el vino]{el agua, el vino — water and wine}

\label{lesson:ES-C26-agua-vino}



\begin{quote}
\textbf{Warm-up.} {[}PAUSE 2s{]} Two drinks, two very different journeys from Latin. One barely changed at all. And the \textquotedblleft{}masculine-looking\textquotedblright{} article on a feminine word isn't a mistake — it's a real, well-documented rule.
\end{quote}

\begin{cousinweb}
\textbf{agua} (\textquotedblleft{}water\textquotedblright{}) $\leftarrow$ Latin \textbf{aqua} — almost unchanged. Compare French, where the very same word was worn all the way down to a bare \textbf{\textquotedblleft{}eau,\textquotedblright{}} just the sound \textquotedblleft{}oh.\textquotedblright{} Spanish kept nearly the whole thing.

English kept the \textbf{original}, whole Latin word too, just as a learned borrowing rather than an everyday one: \textbf{aquatic, aquarium, aqueduct, aqualung} all come straight from \emph{aqua}.
\end{cousinweb}

\begin{grammarlens}[title={Grammar Lens: el agua's article and gender}]
Notice: \textbf{el agua}, not \emph{la agua} — but \emph{agua} is genuinely \textbf{feminine} (you'll see it in \emph{el agua fría}, \textquotedblleft{}the cold water,\textquotedblright{} where \emph{fría} keeps its feminine ending). This isn't an exception to gender; it's a \textbf{pronunciation rule}: Spanish swaps \emph{la} for \emph{el} right before any \textbf{singular} feminine noun that \textbf{starts with a stressed \emph{a}-sound} (to avoid running two identical vowel sounds together, \emph{la agua}) — the plural goes right back to \emph{las}: \emph{las aguas}, never \emph{los aguas}. You'll meet the singular swap again with words like \emph{el águila} (\textquotedblleft{}the eagle,\textquotedblright{} still feminine) and \emph{el hacha} (\textquotedblleft{}the axe,\textquotedblright{} still feminine). It's only the article that changes, never the noun's actual gender.
\end{grammarlens}

\begin{cousinweb}
\textbf{vino} (\textquotedblleft{}wine\textquotedblright{}) $\leftarrow$ Latin \textbf{vīnum} — held its shape, the way French \emph{vin} did too. English shares the very same root, by more than one road: \textbf{wine} itself is \emph{vīnum} borrowed directly into old Germanic; \textbf{vine} and \textbf{vinegar} arrived centuries later, secondhand, via \textbf{French} (\emph{vin + aigre}, \textquotedblleft{}sour wine,\textquotedblright{} for vinegar); and \textbf{vineyard} isn't a borrowing at all — English built it at home, from the already-borrowed \textquotedblleft{}wine\textquotedblright{} plus native \textquotedblleft{}yard.\textquotedblright{}
\end{cousinweb}

\subsection*{Guided Practice}
{[}PAUSE 1s{]}

\begin{itemize}
\raggedright
  \item {[}YOU SAY: \textquotedblleft{}el agua\textquotedblright{} — water, feminine, but takes el{]}
  \item {[}YOU SAY: \textquotedblleft{}el vino\textquotedblright{} — wine{]}
  \item {[}YOU SAY: the contrast — \textquotedblleft{}agua barely changed from aqua; French's eau wore all the way down\textquotedblright{}{]}
  \item {[}YOU SAY: \textquotedblleft{}el águila, el hacha\textquotedblright{} — the same el-before-stressed-a pattern{]}
\end{itemize}

\begin{tcolorbox}[breakable,colback=teal!4,colframe=teal!35!black,title={Wrap-up recall}]
{[}PAUSE 3s{]} Is \textbf{agua} masculine or feminine, and why does it take \emph{el}? (\textbf{Feminine} — \emph{el} replaces \emph{la} only because \emph{agua} starts with a \textbf{stressed a}, a pronunciation rule, not a gender change.) How does \emph{agua}'s erosion from Latin \emph{aqua} compare to French's \emph{eau}? (\textbf{Much less} — Spanish kept nearly the whole word; French wore it down to a bare \textquotedblleft{}oh.\textquotedblright{}) Give an English cousin of \emph{vino}. (\textbf{Wine, vine, vinegar, vintage}.)
\end{tcolorbox}
\section[el pan]{el pan — bread, and the friend who shares it}

\label{lesson:ES-C26-pan}



\begin{quote}
\textbf{Warm-up.} {[}PAUSE 2s{]} Spanish's word for bread hides the same lovely word-story you'd find in English's \textquotedblleft{}companion\textquotedblright{} — except Spanish kept \textbf{both halves}, the bread and the friendship.
\end{quote}

\begin{cousinweb}
\textbf{el pan} (\textquotedblleft{}bread\textquotedblright{}) $\leftarrow$ Latin \textbf{pānis}. Masculine: \emph{el pan}.
\end{cousinweb}

\begin{culture}
The Latin root \textbf{pān-} (\textquotedblleft{}bread\textquotedblright{}) gives Spanish its own word for a close friend:

\begin{itemize}
\raggedright
  \item \textbf{compañero / compañera} (\textquotedblleft{}companion, friend, partner\textquotedblright{}) $\leftarrow$ \emph{com-} (\textquotedblleft{}with\textquotedblright{}) + \emph{pānis} — literally \textquotedblleft{}\textbf{one you share bread with}.\textquotedblright{} Spanish built this word from the very same recipe as English \textbf{companion} and \textbf{company} (which come from the identical Latin pieces) — but where English only kept the \emph{friendship} sense, Spanish still visibly uses the \textbf{same root} for the actual \textbf{bread} (\emph{pan}) too.
  \item \textbf{panadería} (\textquotedblleft{}bakery\textquotedblright{}) and \textbf{panadero} (\textquotedblleft{}baker\textquotedblright{}) — both transparently built on \emph{pan}.
\end{itemize}

So \emph{pan} and \emph{compañero} are the same word-family doing double duty in Spanish: name the food, and name the friend you'd share it with.
\end{culture}

\subsection*{Guided Practice}
{[}PAUSE 1s{]}

\begin{itemize}
\raggedright
  \item {[}YOU SAY: \textquotedblleft{}el pan\textquotedblright{} — bread{]}
  \item {[}YOU SAY: \textquotedblleft{}compañero\textquotedblright{} — literally, one you share bread with{]}
  \item {[}YOU SAY: \textquotedblleft{}panadería, panadero\textquotedblright{} — bakery, baker{]}
\end{itemize}

\begin{tcolorbox}[breakable,colback=teal!4,colframe=teal!35!black,title={Wrap-up recall}]
{[}PAUSE 3s{]} Give \textquotedblleft{}bread\textquotedblright{} with its article. (\textbf{El pan}.) What does \emph{compañero} literally mean, and what's it built from? (\textquotedblleft{}\textbf{One you share bread with}\textquotedblright{} — \emph{com-} + \emph{pānis}.) Does English keep this same double meaning? (\textbf{No} — English's \emph{companion} only kept the friendship sense; Spanish still uses the same root for the actual bread, \emph{pan}, too.)
\end{tcolorbox}
